% !TEX program = xelatex
%==============================================================================
%  ACADEMIC ARTICLE  -  two-column scientific article (XeLaTeX)
%  Fonts : TeX Gyre Pagella (serif text) + TeX Gyre Heros (sans headings)
%  Layout: A4, 10 pt, genuine two-column body with a full-width title block
%  Edit  : the ARTICLE FRONT MATTER block (just above \begin{document}) and the
%          body. The style system in between needs no changes.
%==============================================================================
\documentclass[10pt,a4paper,twocolumn]{article}

%------------------------------------------------------------------------------
%  FONTS  (bundled with the app; headings and metadata use \plumsans)
%------------------------------------------------------------------------------
\usepackage{fontspec}
\defaultfontfeatures{Ligatures=TeX}
\setmainfont{TeX Gyre Pagella}            % primary serif
\newfontfamily\plumsans{TeX Gyre Heros}   % complementary sans-serif

%------------------------------------------------------------------------------
%  PAGE LAYOUT  (margins, column gutter)
%------------------------------------------------------------------------------
\usepackage[a4paper,top=24mm,bottom=26mm,left=18mm,right=18mm,%
  headheight=14pt,headsep=9mm,footskip=12mm]{geometry}
\setlength{\columnsep}{8mm}
\setlength{\columnseprule}{0pt}
\setlength{\parindent}{1.1em}
\setlength{\parskip}{0pt}
\tolerance=900
\emergencystretch=2em
\clubpenalty=5000
\widowpenalty=5000

%------------------------------------------------------------------------------
%  COLORS  (restrained editorial palette; hex swatches are editable)
%------------------------------------------------------------------------------
\usepackage[table]{xcolor}
\definecolor{plumink}{HTML}{20232A}      % title and author names
\definecolor{plumaccent}{HTML}{6E2F4F}   % primary accent: headings, rules, abstract
\definecolor{plumteal}{HTML}{2E5E63}     % secondary accent: links, markers
\definecolor{plummuted}{HTML}{5F6470}    % running heads, affiliations, notes
\definecolor{plumtint}{HTML}{F6F2F5}     % abstract background
\definecolor{plumrule}{HTML}{CDBBC6}     % hairlines
\definecolor{plumorcid}{HTML}{3F7D20}    % ORCID marker

%------------------------------------------------------------------------------
%  PACKAGES
%------------------------------------------------------------------------------
\usepackage{microtype}
\usepackage{amsmath,amssymb,amsthm}
\usepackage{graphicx}
\usepackage{booktabs}
\usepackage{enumitem}
\usepackage{titlesec}
\usepackage{fancyhdr}
\usepackage{caption}
\usepackage{hyperref}

%------------------------------------------------------------------------------
%  LINKS
%------------------------------------------------------------------------------
\hypersetup{colorlinks=true,linkcolor=plumteal,citecolor=plumteal,
  urlcolor=plumaccent,filecolor=plumaccent}

%------------------------------------------------------------------------------
%  HEADINGS  (sans-serif, accent color, thin rule under sections)
%------------------------------------------------------------------------------
\titleformat{\section}
  {\plumsans\bfseries\normalsize\color{plumaccent}}{\thesection}{0.7em}
  {\MakeUppercase}[{\color{plumrule}\titlerule[0.5pt]}]
\titleformat{\subsection}
  {\plumsans\bfseries\small\color{plumink}}{\thesubsection}{0.6em}{}
\titleformat{\subsubsection}[runin]
  {\normalfont\itshape\color{plumink}}{\thesubsubsection}{0.5em}{}[.]
\titlespacing*{\section}{0pt}{2ex plus .4ex minus .2ex}{1ex plus .2ex}
\titlespacing*{\subsection}{0pt}{1.5ex plus .3ex minus .2ex}{0.6ex plus .1ex}
\titlespacing*{\subsubsection}{0pt}{1.2ex plus .2ex}{0.6em}

%------------------------------------------------------------------------------
%  RUNNING HEADS AND FOOTERS  (first page: footer only)
%------------------------------------------------------------------------------
\pagestyle{fancy}
\fancyhf{}
\fancyhead[L]{\plumsans\footnotesize\color{plummuted}\runningauthors}
\fancyhead[R]{\plumsans\footnotesize\color{plummuted}\runningtitle}
\fancyfoot[L]{\plumsans\footnotesize\color{plummuted}\journalname}
\fancyfoot[R]{\plumsans\footnotesize\color{plumaccent}\thepage}
\renewcommand{\headrulewidth}{0.4pt}
\renewcommand{\headrule}{{\color{plumrule}\hrule height \headrulewidth width \headwidth}\vskip-\headrulewidth}
\fancypagestyle{plain}{%
  \fancyhf{}%
  \fancyfoot[L]{\plumsans\footnotesize\color{plummuted}\journalname}%
  \fancyfoot[R]{\plumsans\footnotesize\color{plumaccent}\thepage}%
  \renewcommand{\headrulewidth}{0pt}}

%------------------------------------------------------------------------------
%  THEOREMS, DEFINITIONS, PROOFS, EQUATIONS
%------------------------------------------------------------------------------
\newtheoremstyle{plumplain}
  {6pt plus 2pt minus 1pt}{6pt plus 2pt minus 1pt}{\itshape}{0pt}
  {\plumsans\bfseries\color{plumaccent}}{.}{0.6em}{}
\newtheoremstyle{plumdef}
  {6pt plus 2pt minus 1pt}{6pt plus 2pt minus 1pt}{\normalfont}{0pt}
  {\plumsans\bfseries\color{plumaccent}}{.}{0.6em}{}
\theoremstyle{plumplain}
\newtheorem{theorem}{Theorem}[section]
\newtheorem{lemma}[theorem]{Lemma}
\theoremstyle{plumdef}
\newtheorem{definition}[theorem]{Definition}
\renewcommand{\qedsymbol}{\textcolor{plumaccent}{\ensuremath{\blacksquare}}}
\numberwithin{equation}{section}

%------------------------------------------------------------------------------
%  LISTS, FOOTNOTES, FLOATS (captions, tables)
%------------------------------------------------------------------------------
\setlist{itemsep=1.5pt,topsep=3pt,parsep=0pt,partopsep=0pt,leftmargin=1.5em}
\setlist[itemize]{label=\textcolor{plumaccent}{\textbullet}}
\renewcommand{\footnoterule}{\kern-3pt{\color{plumrule}\hrule width 0.35\columnwidth height 0.4pt}\kern 2.6pt}
\DeclareCaptionFont{plumsans}{\plumsans}
\DeclareCaptionFont{plumaccent}{\color{plumaccent}}
\captionsetup{font=small,labelfont={plumsans,bf,plumaccent},labelsep=period,
  justification=justified,singlelinecheck=false,skip=6pt}
\captionsetup[table]{position=top}

%------------------------------------------------------------------------------
%  TITLE BLOCK AND ABSTRACT  (full width above the two columns)
%  Used as  \twocolumn[{\maketitle \begin{abstract}...\end{abstract}}]
%------------------------------------------------------------------------------
\makeatletter
\gdef\plum@affiliations{}
\gdef\plum@correspondence{}
\newcommand{\affiliations}[1]{\gdef\plum@affiliations{#1}}
\newcommand{\affiliation}[2]{\textsuperscript{\plumsans\color{plumteal}#1}\,#2\par}
\newcommand{\correspondence}[2]{\gdef\plum@correspondence{%
  \textsuperscript{\plumsans\color{plumteal}*}\,Corresponding author: #1, \href{mailto:#2}{#2}}}
\newcommand{\orcid}[1]{\href{https://orcid.org/#1}{%
  \textcolor{plumorcid}{\plumsans\bfseries\scriptsize iD}\,%
  {\plumsans\scriptsize\textcolor{plummuted}{#1}}}}
\newcommand{\plumauthor}[3]{\mbox{%
  {\fontsize{12}{14}\selectfont\color{plumink}#1\textsuperscript{\plumsans\color{plumteal}#2}}%
  \enspace\orcid{#3}}\hspace{1.8em plus 0.8em}}
\renewcommand{\maketitle}{%
  \begingroup
  \normalfont\raggedright\setlength{\parindent}{0pt}%
  \noindent{\color{plumaccent}\rule{\textwidth}{2.2pt}}\par\vspace{6pt}%
  \noindent\makebox[\textwidth][s]{\plumsans\small\bfseries\color{plumaccent}%
    \expandafter\MakeUppercase\expandafter{\articletype}\hfill
    \mdseries\color{plummuted}\@date}\par\vspace{10pt}%
  {\fontsize{22}{26}\selectfont\bfseries\color{plumink}\@title\par}\vspace{12pt}%
  \@author\par\vspace{8pt}%
  {\plumsans\footnotesize\color{plummuted}\plum@affiliations\par}%
  {\plumsans\footnotesize\color{plummuted}\plum@correspondence\par}%
  \vspace{12pt}%
  \endgroup
  \thispagestyle{plain}}
\makeatother
\newsavebox{\plumabsbox}
\renewenvironment{abstract}{%
  \par\setlength{\fboxsep}{10pt}%
  \begin{lrbox}{\plumabsbox}%
  \begin{minipage}{\dimexpr\textwidth-2\fboxsep\relax}%
  \small\setlength{\parindent}{0pt}%
  {\plumsans\bfseries\footnotesize\color{plumaccent}%
    \expandafter\MakeUppercase\expandafter{\abstractname}}\par\vspace{3pt}%
  \noindent\ignorespaces}{%
  \end{minipage}\end{lrbox}%
  \noindent{\color{plumaccent}\rule{\textwidth}{1.2pt}}\par\nointerlineskip
  \noindent\colorbox{plumtint}{\usebox{\plumabsbox}}\par\nointerlineskip
  \noindent{\color{plumrule}\rule{\textwidth}{0.4pt}}\par\vspace{8pt}}
\newcommand{\keywords}[1]{\par\medskip
  \noindent{\color{plumrule}\rule{\linewidth}{0.4pt}}\par\smallskip
  \noindent{\plumsans\bfseries\scriptsize\color{plumaccent}KEYWORDS}\quad #1\par}

%==============================================================================
%  ARTICLE FRONT MATTER  -  edit this block
%==============================================================================
\newcommand{\journalname}{Journal of Illustrative Research}   % footer, left
\newcommand{\articletype}{Research Article}                   % label above the title
\newcommand{\runningtitle}{Short Title of the Article}        % running head, right
\newcommand{\runningauthors}{First Author et al.}             % running head, left

\title{Title of the Article\\[0.5em] {\large A Concise Subtitle That Clarifies the Scope}}
\author{%
  \plumauthor{First Author}{1}{0000-0000-0000-0001}%
  \plumauthor{Second Author}{1,2}{0000-0000-0000-0002}%
  \plumauthor{Third Author}{2,*}{0000-0000-0000-0003}}
\affiliations{%
  \affiliation{1}{Department of Mathematics, University Name, City, Country}%
  \affiliation{2}{Institute of Applied Sciences, Institution Name, City, Country}}
\correspondence{Third Author}{third.author@example.org}
\date{\today}
\hypersetup{pdftitle={Title of the Article},
  pdfauthor={First Author; Second Author; Third Author},
  pdfsubject={}, pdfkeywords={keyword one, keyword two, keyword three}}

\begin{document}

% Full-width title block and abstract, then the two-column body.
% Keep the abstract as one paragraph (no blank lines) inside this bracket.
\twocolumn[{%
  \maketitle
  \begin{abstract}
    Replace this paragraph with a concise summary of the problem, the method and
    the main result, in about 150 words. State what was done, what was found and
    why it matters, and avoid citations or abbreviations that are not explained
    within the abstract itself. This template sets the article in two balanced
    columns with a full-width title block, numbered equations, theorems, figures,
    tables and a reference list.
    \keywords{keyword one; keyword two; keyword three; keyword four}
  \end{abstract}
}]

\section{Introduction}\label{sec:intro}
Scientific articles begin by motivating a problem and stating, as plainly as
possible, what the reader will learn~\cite{smith2020,lee2021}. Replace this text
with the motivation, the gap in the literature and the main contribution.%
\footnote{Footnotes are set at the foot of the column in which they are cited;
keep them short.} The main contributions are:
\begin{itemize}
  \item a formulation of the problem (Section~\ref{sec:methods});
  \item a convergence guarantee (Theorem~\ref{thm:main});
  \item an empirical comparison (Table~\ref{tab:results}).
\end{itemize}
The paper is organized as follows. Section~\ref{sec:background} reviews related
work and fixes notation, Section~\ref{sec:methods} describes the method,
Section~\ref{sec:results} reports results, and Sections~\ref{sec:discussion}
and~\ref{sec:conclusion} discuss and conclude.

\section{Background and Related Work}\label{sec:background}
Early approaches treated the problem in isolation~\cite{garcia2019}, while later
work unified it with a broader family of models~\cite{lee2021}. We write
$\lVert\cdot\rVert$ for the Euclidean norm and
$\langle\cdot,\cdot\rangle$ for the inner product on $\mathbb{R}^{d}$.

\begin{definition}[Smoothness]\label{def:smooth}
A differentiable function $f\colon\mathbb{R}^{d}\to\mathbb{R}$ is
\emph{$L$-smooth} if
$\lVert\nabla f(x)-\nabla f(y)\rVert\le L\lVert x-y\rVert$
for all $x,y\in\mathbb{R}^{d}$.
\end{definition}

\section{Methods}\label{sec:methods}
\subsection{Objective}
We fit a parametric model $f(\cdot\,;\theta)$ by minimizing the regularized
empirical risk
\begin{equation}\label{eq:objective}
  J(\theta)=\mathcal{L}(\theta)+\lambda\,\Omega(\theta),
\end{equation}
where the data term and the regularizer are
\begin{align}
  \mathcal{L}(\theta)&=\frac{1}{n}\sum_{i=1}^{n}\ell\bigl(y_i,f(x_i;\theta)\bigr),\label{eq:loss}\\
  \Omega(\theta)&=\tfrac{1}{2}\lVert\theta\rVert^{2}.\label{eq:reg}
\end{align}
Equation~\eqref{eq:objective} combines~\eqref{eq:loss} and~\eqref{eq:reg}
through the weight $\lambda>0$.

\subsection{Optimization}
Our analysis relies on a standard inequality for smooth functions.
\begin{lemma}[Descent lemma]\label{lem:descent}
If $f$ is $L$-smooth (Definition~\ref{def:smooth}), then for all
$x,y\in\mathbb{R}^{d}$
\begin{align*}
  f(y)&\le f(x)+\langle\nabla f(x),y-x\rangle\\
      &\quad+\tfrac{L}{2}\lVert y-x\rVert^{2}.
\end{align*}
\end{lemma}
\begin{proof}
Put $g(t)=f(x+t(y-x))$. Then $g'(t)-g'(0)\le Lt\lVert y-x\rVert^{2}$ by
Cauchy--Schwarz and smoothness. Integrating over $[0,1]$ gives
$g(1)\le g(0)+g'(0)+\tfrac{L}{2}\lVert y-x\rVert^{2}$, which is the claim.
\end{proof}

\section{Results}\label{sec:results}
The main result concerns plain gradient descent on Eq.~\eqref{eq:objective}.
\begin{theorem}[Convergence]\label{thm:main}
Let $J$ be $L$-smooth and bounded below by $J^{\ast}$. Gradient descent with
step size $1/L$, that is $\theta_{k+1}=\theta_k-\tfrac{1}{L}\nabla J(\theta_k)$,
satisfies
\[
  \min_{0\le k<K}\lVert\nabla J(\theta_k)\rVert^{2}
  \le\frac{2L\,\bigl(J(\theta_0)-J^{\ast}\bigr)}{K}.
\]
\end{theorem}
\begin{proof}
Apply Lemma~\ref{lem:descent} with $x=\theta_k$ and $y=\theta_{k+1}$:
\[
  J(\theta_{k+1})\le J(\theta_k)-\tfrac{1}{2L}\lVert\nabla J(\theta_k)\rVert^{2}.
\]
Summing over $k=0,\dots,K-1$ and telescoping gives
$\tfrac{1}{2L}\sum_{k<K}\lVert\nabla J(\theta_k)\rVert^{2}\le J(\theta_0)-J^{\ast}$.
The minimum is at most the average, which proves the claim.
\end{proof}

Table~\ref{tab:results} compares the proposed method with two baselines, and
Figure~\ref{fig:overview} summarizes the behavior over the iterations. The
values shown are illustrative; replace them with your own data.

\begin{table}[t]
  \centering\small\arrayrulecolor{plumaccent}
  \caption{Mean test error over five runs (illustrative values).}\label{tab:results}
  \begin{tabular}{@{}lccc@{}}
    \toprule
    Method & Error (\%) & Time (s) & Params \\
    \midrule
    Baseline          & 12.4        & 1.8 & 0.9M \\
    Method A          & 10.9        & 2.1 & 1.1M \\
    \textbf{Proposed} & \textbf{9.7} & 2.4 & 1.2M \\
    \bottomrule
  \end{tabular}
\end{table}

\begin{figure}[t]
  \centering
  % Replace the placeholder box with: \includegraphics[width=\linewidth]{figures/your-figure}
  % For a full-width figure use a figure* environment with [t] placement.
  \fcolorbox{plumrule}{plumtint}{\parbox[c][3.2cm][c]{\dimexpr\linewidth-2\fboxsep-2\fboxrule\relax}{%
    \centering\plumsans\small\color{plummuted}Figure placeholder}}
  \caption{Gradient norm versus iteration (placeholder).}\label{fig:overview}
\end{figure}

\section{Discussion}\label{sec:discussion}
Interpret the results in light of the question posed in the introduction and
compare them with prior work~\cite{garcia2019}. Theorem~\ref{thm:main} guarantees
convergence to a stationary point only; it says nothing about global optimality.
Discuss limitations, threats to validity and the assumptions on which the
analysis depends.

\section{Conclusion}\label{sec:conclusion}
Summarize the contribution in two or three sentences and outline directions for
future work.

\section*{Acknowledgments}
The authors thank colleagues and reviewers for their helpful comments.

\section*{Funding}
This work was supported by [funding body] under grant number [000000].

\section*{Conflict of Interest}
The authors declare that they have no known competing financial interests or
personal relationships that could have influenced the work reported in this
article.

\section*{Data Availability}
The data and code supporting the findings of this study are available at
\url{https://example.org/repository}.

\begingroup\small
\begin{thebibliography}{9}
\setlength{\itemsep}{2pt}
\bibitem{smith2020} A.~Smith and B.~Jones, \emph{Title of a journal article},
  Journal of Examples \textbf{12} (2020), 45--67.
\bibitem{lee2021} C.~Lee, \emph{Title of a book}, Publisher, City, 2021.
\bibitem{garcia2019} D.~Garc\'{\i}a, E.~Rossi and F.~Chen, \emph{Title of a
  conference paper}, Proc.\ Conf.\ Examples, 2019, pp.~101--110.
\end{thebibliography}
\endgroup

\end{document}